\section{Moving the Floor: A Mitigation Study}
\label{sec:mitigation}

A benchmark that only reports a failure invites the reading that the
failure is intrinsic. The localisation analysis says otherwise: the
detectors lose sensitivity, not box placement, which is the kind of
deficit training data can address. We test that directly, and we design
the test so that a positive result cannot be an artefact of training on
the condition being scored.

\subsection{Design}

For each of the three SARD detectors we fine-tune the existing clean
checkpoint on training tiles augmented with \cfd{low\_light} at
\textbf{severities 1 and 2 only}. Severity 3, the condition where all
three collapse, is held out of training entirely and is only ever seen
at evaluation. Reporting a recovery at a severity the model was trained
on would measure memorisation of a known condition; holding severity 3
out makes the question generalisation along the ladder, which is the
question a deployment actually faces, since the weather on the night of
a callout is not the weather in the training set.

Three further constraints keep the comparison honest. Validation data
stays clean, enforced by the augmentation script, so the operating point
cannot drift with the augmentation. Each mitigated checkpoint has its
threshold re-selected on its own clean validation data by the same
F1-maximising rule used for every other model in the paper, and then
frozen across all 28 conditions; it does not inherit the baseline's
threshold, because that would score a new model at another model's
operating point. And every mitigated model is re-swept over the full
28-condition grid, not only over low light, so that the cost of the
intervention elsewhere is measured rather than assumed.

\subsection{The collapse is addressable}

Table~\ref{tab:mitigation} gives the result. At the held-out severity 3,
recall moves off the floor for all three architectures: from exactly
0.000 to 0.121 for Faster R-CNN, 0.059 for RT-DETR and 0.188 for
YOLOv11, with bootstrap intervals that exclude zero in every case. At
the trained severities the recovery is much larger, from 0.419 to 0.843,
0.260 to 0.772 and 0.184 to 0.807 at severity 1, and from at most 0.077
to between 0.529 and 0.647 at severity 2.

Two readings are available and we take the conservative one. A
severity-3 recall of 0.059 to 0.188 is not an operational fix; a system
that finds one subject in five at late dusk is still failing. What the
result establishes is narrower and sufficient for the benchmark's
purpose: the zero is not a hard limit of these architectures under
photon-limited capture, it is a data-coverage deficit that responds to
targeted training, and it responds even one full severity step beyond
the training distribution. The benchmark therefore measures something
that can be optimised against, which is the property that makes an
evaluation axis useful rather than merely discouraging.

\begin{table}[t]
\centering\small
\setlength{\tabcolsep}{3pt}
\caption{Low-light mitigation on SARD. Each model is fine-tuned on
severities 1 and 2 only; \textbf{severity 3 is held out of training} and
seen only at evaluation. Thresholds are re-selected on clean validation
data for each mitigated checkpoint and frozen. Bootstrap 95\% intervals
over 1000 resamples for the mitigated models.}
\label{tab:mitigation}
\begin{tabular}{llcccc}
\toprule
Model & & Clean & s1 & s2 & s3 \\
\midrule
\multirow{3}{*}{F.\ R-CNN}
 & before & 0.879 & 0.419 & 0.077 & 0.000 \\
 & after  & 0.900 & 0.843 & 0.641 & \textbf{0.121} \\
 & 95\%   &       & \ci{0.818}{0.865} & \ci{0.612}{0.669} & \ci{0.101}{0.140} \\
\midrule
\multirow{3}{*}{RT-DETR}
 & before & 0.888 & 0.260 & 0.013 & 0.000 \\
 & after  & 0.858 & 0.772 & 0.529 & \textbf{0.059} \\
 & 95\%   &       & \ci{0.745}{0.798} & \ci{0.497}{0.560} & \ci{0.046}{0.073} \\
\midrule
\multirow{3}{*}{YOLOv11}
 & before & 0.891 & 0.184 & 0.007 & 0.000 \\
 & after  & 0.883 & 0.807 & 0.647 & \textbf{0.188} \\
 & 95\%   &       & \ci{0.782}{0.830} & \ci{0.615}{0.678} & \ci{0.168}{0.211} \\
\bottomrule
\end{tabular}
\end{table}

\subsection{What it costs, and why the cost is not one number}

The interesting part is the price. Because each mitigated model was
re-swept over the whole grid, we can measure the change on the eight
corruptions that were never trained on, 24 conditions per model, as well
as on clean data (Table~\ref{tab:tradeoff}).

The result does not resolve into a single robustness tax. Faster R-CNN
improves on average by 3.7 points across the other 24 conditions, gains
2.1 points on clean data, and its worst single regression is 2.5 points,
so nothing degrades by more than the noise band of 3 points that we use
as a reporting cutoff. RT-DETR moves the other way: it loses 4.3 points
on average, regresses by more than 3 points on 17 of 24 conditions, loses
3.0 points on clean data, and its worst case is 11.6 points. YOLOv11 sits
between them, losing 1.9 points on average with 8 of 24 conditions
regressing by more than 3 points.

Because all three arms use the same augmented data, the same held-out
severity, the same threshold-selection rule and the same evaluator, the
spread across them is attributable to architecture rather than to
protocol. We do not claim a mechanism. We note that the transformer
detector, which already carries the largest relative recall drop of the
three in every disaster family on SARD (Table~\ref{tab:family}), is also
the one for which targeted augmentation is most
expensive elsewhere, and that this is precisely the kind of interaction
that a benchmark reporting one averaged robustness number could not
surface. The practical implication for a deployment is that
``augment for the condition you expect'' is sound advice for some
architectures and a regression for others, and which one it is has to be
measured on the full grid rather than inferred.

\begin{table}[t]
\centering\small
\setlength{\tabcolsep}{3pt}
\caption{Cost of low-light augmentation on conditions it was not trained
for: change in recall over the eight other corruptions at three
severities (24 conditions per model), plus clean. Positive is better.
The sign of the mean differs by architecture.}
\label{tab:tradeoff}
\begin{tabular}{lcccc}
\toprule
 & $\Delta$ clean & Mean $\Delta$ & Worst & \#\,worse \\
 &                & (24 cond.)    & $\Delta$ & by $>$0.03 \\
\midrule
F.\ R-CNN & $+0.021$ & $+0.037$ & $-0.025$ & 0/24 \\
YOLOv11   & $-0.007$ & $-0.019$ & $-0.066$ & 8/24 \\
RT-DETR   & $-0.030$ & $-0.043$ & $-0.116$ & 17/24 \\
\bottomrule
\end{tabular}
\end{table}

\subsection{Scope of this study}

The mitigation arm is one augmentation strategy, on one dataset, at one
seed per architecture, targeting one corruption. It is enough to settle
whether the floor is movable and to show that the cost of moving it is
architecture-dependent. It is not enough to recommend a training recipe,
and we do not. The natural extensions, sweeping the augmentation across
the other corruptions, repeating it on HERIDAL, and combining
conditions as a real mission would, are supported by the released
pipeline and are left open.
